\subsection{Diagnostic 1D-3D Velocity Comparison}
\label{sec:resolved-3d-comparison}

\subsubsection{Available 3D Data and Matching Limits}
\label{subsec:preliminary-cross-model-velocity-comparison}

This section reports the final-time resolved-velocity comparison whose protocol
is defined in Section~\ref{subsec:numerical-experiments-metrics-provenance}. It
is a cross-model comparison of velocity and physical-flow observables under the
declared plane--tetrahedron quadrature operator. The zero-input rest-state test
in Section~\ref{sec:numerical-verification} remains required context because it
uses the same principal operator as the comparison runs.

The available resolved datasets are the two local velocity fields labelled C23
and C40. The comparison record uses source ID 77 for C23 and source ID 60 for
C40, each with an XDMF velocity sample at $0.9995$ s, while the corresponding
1D run completes at $1.0$ s. The $5\times10^{-4}$ s sample-time offset is
reported, not eliminated. The current/deformed 3D geometry status also remains
unresolved: the tracked assets support the static cut-area audit below, but they
do not establish whether a deformed-wall geometry or displacement application
should be used for the final-time plane cuts. The 3D wall model, wall-motion
history, inlet/outlet histories, material parameters, and prior transient
history are likewise not persisted in the tracked report artifact.

\begin{table}[!htb]
    \centering
    \scriptsize
    \caption{Recorded subcritical-boundary diagnostics for the same C23/C40
    1D runs used in the resolved-velocity comparison. The positive
    characteristic-radicand values are in solver units of cm$^2$/s$^2$; the
    persisted $\lambda_-$, $\lambda_+$, and subcriticality margin support the
    boundary-regime gate in Assumption~\ref{ass:1d-subcritical-boundary-regime}.}
    \label{tab:t1-alpha-radicand-diagnostics}
    \resizebox{\textwidth}{!}{%
    \begin{tabular}{@{}lrrrrrrrr@{}}
        \toprule
        \textbf{Case}
            & $\boldsymbol{\min \alpha_{\mathrm{eff}}}$
            & $\boldsymbol{\max \alpha_{\mathrm{eff}}}$
            & $\boldsymbol{\min \chi^2}$
            & $\boldsymbol{\min\lambda_-}$
            & $\boldsymbol{\max\lambda_-}$
            & $\boldsymbol{\min\lambda_+}$
            & $\boldsymbol{\max\lambda_+}$
            & $\boldsymbol{\min\{\lambda_+,-\lambda_-\}}$ \\
        \midrule
        C23 & 1.33058 & 1.33333 & $8.15\times 10^5$ & -998.85 & -852.17 & 953.33 & 1058.76 & 852.17 \\
        C40 & 1.32500 & 1.33333 & $6.36\times 10^5$ & -999.08 & -713.79 & 880.69 & 1058.94 & 713.79 \\
        \bottomrule
    \end{tabular}
    }
\end{table}

The entries in Table~\ref{tab:t1-alpha-radicand-diagnostics} are used as the
subcritical-boundary gate for the paired C23/C40 comparison runs; they do not
widen the velocity comparison beyond the declared observation operator.

\FloatBarrier

\subsubsection{Observation Operator}
\label{subsec:cross-sectional-sampling-protocol-results}

The comparison uses the declared
\texttt{CrossSectionQuadratureOperator}. For each target plane, the operator
intersects the tetrahedral 3D mesh with the plane, triangulates the cut polygon,
and computes section area, physical flow, and area-mean axial velocity from the
cut triangles as stated in Section~\ref{subsec:plane-tetrahedron-comparison-operator}.
The 1D quantities use the solver-to-physical map from
Definition~\ref{def:physical-to-solver-map}: $Q_{1D}=\pi q$ and
$\bar u_{1D}=q/a$. The principal comparison uses 200 equally spaced section
targets on $0\le z\le6$ cm. Secondary radial-profile outputs were generated at
$z=1.951$, 2.451, and 2.951 cm, but they are quarantined rather than used as
retained results.

The stenosis throat is the minimizer $z^\star=\argmin_z R_0(z)$ of the smooth
reference-radius profile in Eq.~\ref{eq:1d-stenosis-profile}. For the comparison
geometry this gives $z^\star\approx2.451\,\mathrm{cm}$ for both severities, with
$R_{\min}=0.1386\,\mathrm{cm}$ for C23 and $R_{\min}=0.1080\,\mathrm{cm}$ for
C40. The shaded region in Figure~\ref{fig:t1-section-mean-comparison} is the
displayed interval $2.4\le z\le2.8\,\mathrm{cm}$.

\subsubsection{Cut-Area Audit}
\label{subsec:section-averaged-velocity-results}

The same quadrature record audits whether the cut operator is sampling the
expected static cross-section. With the reference area
$A_{\mathrm{ref}}(z_j)=\pi R_0(z_j)^2$, define
\begin{flalign*}
    \quad
    \epsilon_A(z_j)
    =
    \frac{|A_{3D}(z_j)-A_{\mathrm{ref}}(z_j)|}
         {A_{\mathrm{ref}}(z_j)}. &&
\end{flalign*}
Table~\ref{tab:t1-area-audit} reports this geometry/operator audit for the same
200 section targets. It shows sub-percent typical static cut-area discrepancy,
with the largest differences confined to a few targets. This audit does not
resolve the current/deformed geometry question and does not measure the dynamic
area discrepancy entering the 1D velocity denominator,
\begin{flalign*}
    \quad
    \epsilon_A^{\mathrm{model}}(z_j,t)
    =
    \frac{|\pi a(z_j,t)-A_{3D}(z_j)|}{A_{3D}(z_j)}. &&
\end{flalign*}

\begin{table}[!htb]
    \centering
    \scriptsize
    \caption{Quadrature cut-area audit against the static reference
    $A_{\mathrm{ref}}(z)=\pi R_0(z)^2$. Discrepancies are percentages over
    valid section targets and are stored in the tracked
    \texttt{area-audit.dat} report asset.}
    \label{tab:t1-area-audit}
    \begin{tabular}{@{}lrrrrr@{}}
        \toprule
        \textbf{Case}
            & \textbf{Sections}
            & $\boldsymbol{\min \epsilon_A}$
            & $\boldsymbol{\mathrm{median}\ \epsilon_A}$
            & $\boldsymbol{\mathrm{mean}\ \epsilon_A}$
            & $\boldsymbol{\max \epsilon_A}$ \\
        \midrule
        C23 & 200 & 0.045 & 0.281 & 0.389 & 2.94 \\
        C40 & 200 & 0.005 & 0.291 & 0.397 & 4.46 \\
        \bottomrule
    \end{tabular}
\end{table}

\subsubsection{Axial Physical-Flow Behavior}
\label{subsec:axial-flow-behavior}

The reported section data include physical flows $Q_{1D}(z)=\pi q(z)$ and
$Q_{3D}(z)$ before they are divided by area to form mean velocity. The axial 3D
flow extracted by the declared plane operator is not constant across the 200
section targets: its range is 7.7\% of the C23 mean and 20.2\% of the C40 mean.
That unresolved axial variation appears before the velocity interpretation
because it is a matching gate, not a solved mechanism. The tracked artifacts do
not identify whether the variation comes from 3D numerical state, boundary or
history mismatch, geometry handling, operator sensitivity, or another source.

The section means are
$\langle Q_{1D}\rangle=2.288\,\mathrm{cm}^3/\mathrm{s}$ and
$\langle Q_{3D}\rangle=2.241\,\mathrm{cm}^3/\mathrm{s}$ for C23, and
$\langle Q_{1D}\rangle=2.287\,\mathrm{cm}^3/\mathrm{s}$ and
$\langle Q_{3D}\rangle=2.219\,\mathrm{cm}^3/\mathrm{s}$ for C40. The signed mean
physical-flow biases are 0.0472 and 0.0684 cm$^3$/s, while the mean absolute
physical-flow discrepancies are 0.0482 and 0.0709 cm$^3$/s and the RMS
physical-flow discrepancies are 0.0509 and 0.0893 cm$^3$/s. The largest
physical-flow discrepancies in the retained section table occur at the same
locations as the largest mean-velocity discrepancies.

\begin{figure}[!htb]
    \centering
    \captionsetup{aboveskip=2pt,belowskip=2pt}
    \figtikz{axial-flow-comparison}
    \caption{Axial physical-flow comparison under the declared plane-quadrature
    operator. Dashed curves are $Q_{1D}(z)=\pi q(z)$ and solid curves are
    $Q_{3D}(z)$ from the available resolved velocity fields. The variation in
    the solid curves is a matching limit for interpreting the velocity
    discrepancies.}
    \label{fig:t1-axial-flow-comparison}
\end{figure}

The algebraic split used for flow/area bookkeeping is
\begin{flalign*}
    \quad
    \frac{Q_{1D}}{A_{1D}}-\frac{Q_{3D}}{A_{3D}}
    =
    \frac{Q_{1D}-Q_{3D}}{A_{1D}}
    +
    Q_{3D}\left(\frac{1}{A_{1D}}-\frac{1}{A_{3D}}\right),
    \qquad A_{1D}=\pi a. &&
\end{flalign*}
The tracked section-quadrature data store enough section-wise values to
reconstruct this split on retained sections with nonzero mean-velocity
denominators: $A_{1D}=Q_{1D}/\bar u_{1D}$ and
$A_{3D}=Q_{3D}/\bar u_{3D}$. Table~\ref{tab:t1-flow-area-decomposition}
summarizes the resulting contribution magnitudes. This decomposition is
bookkeeping for the reported samples, not causal attribution.

\begin{table}[!htb]
    \centering
    \scriptsize
    \caption{Section-wise decomposition of the mean-velocity discrepancy into
    physical-flow and area-denominator contributions. Values are in cm/s and are
    reconstructed from the tracked section-quadrature report asset. The split is
    bookkeeping, not a causal attribution.}
    \label{tab:t1-flow-area-decomposition}
    \resizebox{\textwidth}{!}{%
    \begin{tabular}{@{}lrrrrl@{}}
        \toprule
        \textbf{Case}
            & $\boldsymbol{\langle |(Q_{1D}-Q_{3D})/A_{1D}|\rangle}$
            & $\boldsymbol{\langle |Q_{3D}(A_{1D}^{-1}-A_{3D}^{-1})|\rangle}$
            & $\boldsymbol{z_{\max |d_u|}}$
            & $\boldsymbol{|d_u|_{\max}}$
            & \textbf{Larger bookkeeping term at max} \\
        \midrule
        C23 & 0.529 & 0.112 & 2.291 & 2.35 & Flow term \\
        C40 & 1.012 & 0.133 & 2.261 & 9.92 & Flow term \\
        \bottomrule
    \end{tabular}%
    }
\end{table}

\subsubsection{Section-Mean Velocity Discrepancies}
\label{subsec:section-mean-velocity-discrepancies}

Table~\ref{tab:t1-3d-comparison-summary} reports the retained section-mean
velocity discrepancy measures. The signed mean velocity biases are 0.480 cm/s
for C23 and 0.909 cm/s for C40, while the mean absolute velocity discrepancies
are 0.505 and 0.966 cm/s. The RMS velocity discrepancy is 0.664 cm/s for C23
and 1.87 cm/s for C40, with relative RMS values 0.0277 and 0.0688. The maximum
section discrepancies are localized in the upstream shoulder: 2.35 cm/s at
$z=2.291$ cm for C23 and 9.92 cm/s at $z=2.261$ cm for C40.

\begin{table}[!htb]
    \centering
    \scriptsize
    \caption{Quadrature-backed section-mean velocity discrepancy summary at
    nominally aligned final-time samples. These are single-realization
    descriptors in cm/s except for the final axial-location column in cm. Angle
    brackets denote axial sample means over the 200 retained comparison
    sections, not temporal averages or inner products.}
    \label{tab:t1-3d-comparison-summary}
    {\setlength{\tabcolsep}{1.8pt}
    \resizebox{\textwidth}{!}{%
    \begin{tabular}{@{}lrrrrrrrrr@{}}
        \toprule
        \textbf{Case}
            & \textbf{Sections}
            & $\boldsymbol{\langle\bar u_{1D}\rangle}$
            & $\boldsymbol{\langle\bar u_{3D}\rangle}$
            & $\boldsymbol{\overline d_u}$
            & $\boldsymbol{D_{u,1}}$
            & $\boldsymbol{D_{u,2}}$
            & \textbf{rel.} $\boldsymbol{D_{u,2}}$
            & $\boldsymbol{\max |d_u|}$
            & $\boldsymbol{z_{\max |d_u|}}$ \\
        \midrule
        C23 & 200 & 24.11 & 23.63 & 0.480 & 0.505 & 0.664 & 0.0277 & 2.35 & 2.291 \\
        C40 & 200 & 26.43 & 25.52 & 0.909 & 0.966 & 1.87 & 0.0688 & 9.92 & 2.261 \\
        \bottomrule
    \end{tabular}
    }%
    }
\end{table}

Figure~\ref{fig:t1-section-mean-comparison} shows the same area-mean axial
velocity comparison along the vessel. The 1D and quadrature-averaged 3D means
are closest away from the largest acceleration zone, while the largest
discrepancies in this two-case pair are upstream of the displayed throat-centered
band. The largest C40 velocity discrepancy occurs at the same section as the
largest static cut-area discrepancy in Table~\ref{tab:t1-area-audit}. Excluding
the two C40 sections with static area discrepancy above 4\% moves the maximum
from 9.92 cm/s at $z=2.261$ cm to 8.74 cm/s at $z=2.291$ cm, where the static
area discrepancy is below 1\%. Thus the localized C40 peak is not solely an
area-audit artifact, but the exact maximum value and location remain
operator-sensitive.

\begin{figure}[H]
    \centering
    \captionsetup{aboveskip=2pt,belowskip=2pt}
    \figtikz{section-mean-comparison}
    \caption{Area-mean axial velocity comparison at $t=1.0$ s. Solid curves are
    tetrahedral cross-section quadrature means from the resolved 3D velocity
    field; dashed curves are interpolated 1D means
    $\bar u=q/a=Q_{\mathrm{phys}}/A_{\mathrm{phys}}$. The shaded band marks the
    displayed throat-centered stenosis region; the largest listed section
    discrepancies occur just upstream of this band.}
    \label{fig:t1-section-mean-comparison}
\end{figure}

\begin{table}[!htb]
    \centering
    \scriptsize
    \caption{Five largest section-wise velocity discrepancies across the two
    included cases. Velocities and velocity discrepancies are in cm/s; physical
    flows and flow discrepancies are in cm$^3$/s, with $Q_{1D}=\pi q$.}
    \label{tab:t1-largest-section-discrepancies}
    \begin{tabular}{@{}rlrrrrrrrr@{}}
        \toprule
        \textbf{Rank}
            & \textbf{Case}
            & $\boldsymbol{z}$ \textbf{(cm)}
            & $\boldsymbol{\bar u_{1D}}$
            & $\boldsymbol{\bar u_{3D}}$
            & $\boldsymbol{|d_u|}$
            & $\boldsymbol{Q_{1D}}$
            & $\boldsymbol{Q_{3D}}$
            & $\boldsymbol{|d_Q|}$
            & \textbf{Tri.} \\
        \midrule
        1 & C40 & 2.261 & 52.58 & 42.66 & 9.92 & 2.201 & 1.862 & 0.339 & 759 \\
        2 & C40 & 2.291 & 56.94 & 48.20 & 8.74 & 2.230 & 1.903 & 0.327 & 749 \\
        3 & C40 & 2.352 & 61.43 & 53.22 & 8.21 & 2.270 & 1.980 & 0.290 & 743 \\
        4 & C40 & 2.322 & 59.81 & 52.09 & 7.72 & 2.253 & 1.970 & 0.283 & 771 \\
        5 & C40 & 2.231 & 46.95 & 39.94 & 7.01 & 2.174 & 1.918 & 0.256 & 833 \\
        \bottomrule
    \end{tabular}
\end{table}

\subsubsection{Quarantined Radial Output}
\label{subsec:reconstructed-radial-profile-results}

The radial-profile reducer is excluded from the retained publication results.
Direct equal-bin recomputation from the tracked radial data has not been
reconciled with the existing radial-summary values. Therefore no radial figure,
table, discrepancy metric, profile-behavior statement, or radial-bin
localization claim is used as evidence in this chapter.

\subsubsection{Interpretation Limits}
\label{subsec:spatial-localization-discrepancies}

The result of this chapter is a localization of descriptive velocity
discrepancies under the declared observation operator. The section-mean
comparison localizes the largest C40 discrepancy upstream of the displayed
throat-centered band. This observation is bounded by the matching limits stated
at the start of the chapter: unmatched or unpersisted 3D wall, boundary,
material, history, and current/deformed geometry information; unresolved axial
variation in extracted 3D flow; and the $1.0$ s versus $0.9995$ s sample-time
offset. The flow/area decomposition is bookkeeping, and the displayed
differences do not identify a unique mechanism.
